\section{Background}
Mechanistic interpretability currently can do supervised steering well
or unsupervised steering poorly. This has prompted the Google Deep Mind
interpretability team to pivot to a pragmatic steering
agenda\cite{gdmprogress2025}. Compared to unsupervised interpretability,
supervised steering is highly robust. However it can only distinguish
what the experimenter is looking for and does not account for the actual
computations performed by a model\cite{makelov2023}. The dominant
unsupervised paradigm is the sparse autoencoder (SAE). A model's hidden
state is approximated as a finite state machine consisting of discrete
features. The initial implementation used a $\relu$ activation and
sparsity loss\cite{bricken2023}, though more recent implementations have
used a $\topk$ activation function without a sparsity loss\cite{gao2024}.

\subsection{Common problems with SAEs}
The motivation is to decompile a model to a finite state machine (FSM).
The limitation is that such an FSM is only as interpretable as an
enumeration of all possible states.

\subsubsection{Features are not composable}
A persistent problem with SAEs is feature absorption\cite{chanin2025}.
Instead of representing an input as a composition of monosemantic
features, an SAE will often instead use a single separate feature. This
is a problem for interpretability, as we would like features to
generalize. One attempt to address this problem is Matryoshka
SAEs\cite{bussmann2025}, which use multiple layers to build a
hierarchical coarse-to-fine representation.

\subsubsection{Features are not causal}
Sparsity trades off against reconstruction fidelity. Using features for
steering often results in inconsistent, context-sensitive
behavior\cite{duan2026}.

\subsubsection{Features are not closed-form}
A major concern is how well an interpretability technique generalizes
out-of-distribution. Ideally we would like to infer behavior by looking
solely at model weights, which can be achieved with bilinear
MLPs\cite{pearce2025}.

\subsubsection{Forward passes are expensive}
Efficiency is also a problem. Although the sparsity constraint means
that only a handful of features are active for a given input, the
entire feature space needs to be computed for each sample.
Mixture-of-experts (MoE) SAE variants reduce compute costs by dividing
feature space between smaller submodules. Examples include Switch
SAEs\cite{mudide2025} and Tree SAEs\cite{cao2026}.

Less explored is the effect of the structure imposed by the MoEfied
architecture on the interpretability of feature space. A common approach
to scalably interpreting features is using an LLM to predict activations
from examples\cite{paulo2025}. With an unstructured feature space, it
quickly becomes intractable to interpret all permutations of features
this way. It is however possible to instead apply this technique to
expert routings. The MLPs of a transformer can themselves be converted
to an MoE architecture\cite{zhang2022}. Language models seem to organize
contrastive concepts into orthogonal simplices\cite{park2025}. This is
suggestive of MoEfication revealing latent contrastive steering vectors.
Imposing an orthogonality constraint has been shown to significantly
reduce feature absorption\cite{korznikov2025}.

Existing attempts to route activations to expert SAEs focus on routing
similar inputs to the same expert. Contrastive learning wants the
opposite: training on things at opposite ends of a shared axis. We don't
just want to find similar features, we want to find contrastive vectors
in feature space that correspond to what a model considers a natural
category.


